\lyseis{Β2.5}{Κέντρο συμμετρίας σχήματος}
\sol{1} Α: όχι. Η: ναι. Ι: ναι. Λ: όχι. Ν: ναι. Ο: ναι. Σ: ναι. Χ: ναι.
\sol{2} Ισοσκελές τρίγωνο: όχι. Ορθογώνιο: ναι (η τομή των διαγωνίων).
\sol{3} Άπειρους άξονες (κάθε διάμετρος) και ένα κέντρο συμμετρίας (το κέντρο του).
\sol{4} Το κέντρο είναι η τομή των διαγωνίων. Οι μισές διαγώνιοι είναι $3$ cm και $4$ cm, άρα δύο κορυφές απέχουν $3$ cm και οι άλλες δύο $4$ cm από το κέντρο. Η πλευρά του ρόμβου είναι $\sqrt{3^2+4^2}=5$ cm.
\sol{5} Ελεύθερη κατασκευή: το σημείο τομής των διαγωνίων.
\sol{6} Ελεύθερη κατασκευή: προκύπτει παράλληλη ευθεία.
\sol{7} α) Λάθος (ισοσκελές τρίγωνο). β) Σωστό. γ) Σωστό.
\sol{8} Σημαίνει ότι έχει κέντρο συμμετρίας: το σημείο τομής των διαγωνίων.
\sol{9} Στην τομή των διαγωνίων: $15$ cm από τις δύο μικρές πλευρές και $10$ cm από τις δύο μεγάλες.
\sol{10} $4$ ζεύγη.
\sol{11} STOP (οκτάγωνο): ναι. Η «Απαγορεύεται η είσοδος» (κύκλος με οριζόντια λωρίδα): ναι.
\sol{12} Ελεύθερη κατασκευή: το μέσο της κεντρικής ευθείας.
\sol{13} Παραλληλόγραμμο (όχι ορθογώνιο ή ρόμβος).
\sol{14} Ισοσκελές τρίγωνο.
\sol{15} Ναι: στην τομή των αξόνων.
\sol{16} Αν $\mathrm{ΟΑ}\perp\varepsilon$, τότε η $\varepsilon'$ είναι κάθετη στην $\mathrm{ΑΑ}'$ στο $\mathrm{Α}'$. Και οι δύο κάθετες στην ίδια ευθεία $\mathrm{ΑΑ}'$ είναι παράλληλες.
\sol{17} Ελεύθερη απάντηση (π.χ. ρολόι, τραπουλόχαρτο, φαγητό).
\sol{18} Ελεύθερη κατασκευή.
